% from ega4-8-fr.tex
\oldpage[IV-3]{27}
assez grand, on voit finalement qu'on est ramené (en remplaçant $S_\lambda$
par $U_\lambda$) au cas où $Z'_\lambda$ et $Z''_\lambda$ sont des
sous-préschémas fermés de $S_\lambda$. Mais alors, cela a déjà été démontré
puisque (8.6.2.3) est croissante et injective.
\end{proof}

\begin{corollary}[8.6.4]
\label{IV.8.6.4-fr}
Supposons $S_0$ quasi-compact, et soit $Z_\lambda$ un sous-préschéma de
$S_\lambda$, de présentation finie sur $S_\lambda$. Pour que
$Z=u_\lambda^{-1}(Z_\lambda)$ soit un sous-préschéma induit sur un ouvert
(resp. un sous-préschéma fermé) de $S$, il faut et il suffit qu'il existe
$\mu\geq\lambda$ tel que $Z_\mu=u_{\lambda\mu}^{-1}(Z_\lambda)$ soit induit
sur un ouvert (resp. un sous-préschéma fermé) dans $S_\mu$.
\end{corollary}

\begin{proof}
Soit $(U_\lambda^{(i)})$ un recouvrement ouvert affine fini de $S_\lambda$, et
posons $U_\mu^{(i)}=u_{\lambda\mu}^{-1}(U_\lambda^{(i)})$ pour
$\mu\geq\lambda$ et $U^{(i)}=u_\lambda^{-1}(U_\lambda^{(i)})$. Si $Z$ est
ouvert (resp. fermé) dans $S$, $Z\cap U^{(i)}$ l'est dans $U^{(i)}$, et
inversement si chacun des $Z_\mu\cap U_\mu^{(i)}$ est ouvert (resp. fermé)
dans $U_\mu^{(i)}$, $Z_\mu$ l'est dans $S_\mu$. Puisque $L$ est filtrant, il
suffit de démontrer le corollaire lorsque $S_\lambda$ est affine, donc
quasi-séparé. Mais alors le résultat découle de ce que les applications
(8.6.2.1), (8.6.2.2) et (8.6.2.3) sont bijectives.
\end{proof}

\subsection{Critères pour qu'une limite projective de préschémas soit un préschéma réduit (resp. intègre)}
\label{subsection:IV.8.7-fr}

Nous conservons les hypothèses et notations de (8.2.2) et supposons toujours
que $S_0$ est un des $S_\lambda$.

\begin{proposition}[8.7.1]
\label{IV.8.7.1-fr}
Supposons que $S$ soit non réduit. Alors il existe $\lambda_0$ tel que pour
$\lambda\geq\lambda_0$, $S_\lambda$ soit non réduit.
\end{proposition}

\begin{proof}
La question étant locale sur $S_0$, on peut supposer
$S_0=\operatorname{Spec}(A_0)$ affine, d'où $S=\operatorname{Spec}(A)$, où
$A=\varinjlim A_\lambda$ est la limite inductive d'un système inductif de
$A_0$-algèbres $(A_\lambda)$. On sait alors (5.13.2) que le nilradical de $A$
est la limite inductive de ceux des $A_\lambda$~; s'il n'est pas nul, un des
$A_\lambda$ contient donc un élément nilpotent $a_\lambda\neq0$ dont l'image
dans $A$ est un élément nilpotent et $\neq0$, et l'image de $a_\lambda$ dans
les $A_\mu$ pour $\mu\geq\lambda$ est par suite un élément nilpotent et
$\neq0$.
\end{proof}

\begin{proposition}[8.7.2]
\label{IV.8.7.2-fr}
Supposons vérifiée l'une des hypothèses suivantes~:
\begin{enumerate}
\item[\emph{a)}] $S_0$ est quasi-compact, le nilradical $\mathcal N_0$ de
$\mathcal O_{S_0}$ est un Idéal de type fini (ce qui sera par exemple vérifié
lorsque $S_0$ est noethérien), et les morphismes
$u_{\lambda\mu}:S_\mu\to S_\lambda$ sont des immersions ouvertes.
\item[\emph{b)}] Les morphismes $u_{\lambda\mu}:S_\mu\to S_\lambda$ sont
fidèlement plats.
\end{enumerate}
Sous ces conditions, pour que $S$ soit réduit, il faut et il suffit qu'il
existe $\lambda_0$ tel que $S_\lambda$ soit réduit pour
$\lambda\geq\lambda_0$.

En outre, dans le cas \emph{b)}, si $S$ est réduit, tous les $S_\lambda$ le
sont.
\end{proposition}

\begin{proof}
La dernière assertion résulte de ce que le morphisme $S\to S_\lambda$ est
alors fidèlement plat pour tout $\lambda$ (8.3.8) et de (2.1.13). D'autre part,
(8.7.1) prouve que la condition de l'énoncé est suffisante (sans hypothèse sur
$S_0$ ni sur les $u_{\lambda\mu}$). Il reste donc à voir que la condition est
nécessaire dans l'hypothèse \emph{a)}~; alors (8.2.13), l'espace sous-jacent à
$S$ s'identifie à l'intersection des espaces sous-jacents aux $S_\lambda$ (les
$S_\lambda$ étant identifiés à des sous-préschémas induits sur des ouverts de
$S_0$), et le faisceau structural $\mathcal O_S$ s'identifie au
